% Options for packages loaded elsewhere
\PassOptionsToPackage{unicode}{hyperref}
\PassOptionsToPackage{hyphens}{url}
\PassOptionsToPackage{dvipsnames,svgnames,x11names}{xcolor}
%
\documentclass[
  12pt]{article}

\usepackage{amsmath,amssymb}
\usepackage{setspace}
\usepackage{iftex}
\ifPDFTeX
  \usepackage[T1]{fontenc}
  \usepackage[utf8]{inputenc}
  \usepackage{textcomp} % provide euro and other symbols
\else % if luatex or xetex
  \usepackage{unicode-math}
  \defaultfontfeatures{Scale=MatchLowercase}
  \defaultfontfeatures[\rmfamily]{Ligatures=TeX,Scale=1}
\fi
\usepackage{lmodern}
\ifPDFTeX\else  
    % xetex/luatex font selection
\fi
% Use upquote if available, for straight quotes in verbatim environments
\IfFileExists{upquote.sty}{\usepackage{upquote}}{}
\IfFileExists{microtype.sty}{% use microtype if available
  \usepackage[]{microtype}
  \UseMicrotypeSet[protrusion]{basicmath} % disable protrusion for tt fonts
}{}
\makeatletter
\@ifundefined{KOMAClassName}{% if non-KOMA class
  \IfFileExists{parskip.sty}{%
    \usepackage{parskip}
  }{% else
    \setlength{\parindent}{0pt}
    \setlength{\parskip}{6pt plus 2pt minus 1pt}}
}{% if KOMA class
  \KOMAoptions{parskip=half}}
\makeatother
\usepackage{xcolor}
\setlength{\emergencystretch}{3em} % prevent overfull lines
\setcounter{secnumdepth}{5}
% Make \paragraph and \subparagraph free-standing
\makeatletter
\ifx\paragraph\undefined\else
  \let\oldparagraph\paragraph
  \renewcommand{\paragraph}{
    \@ifstar
      \xxxParagraphStar
      \xxxParagraphNoStar
  }
  \newcommand{\xxxParagraphStar}[1]{\oldparagraph*{#1}\mbox{}}
  \newcommand{\xxxParagraphNoStar}[1]{\oldparagraph{#1}\mbox{}}
\fi
\ifx\subparagraph\undefined\else
  \let\oldsubparagraph\subparagraph
  \renewcommand{\subparagraph}{
    \@ifstar
      \xxxSubParagraphStar
      \xxxSubParagraphNoStar
  }
  \newcommand{\xxxSubParagraphStar}[1]{\oldsubparagraph*{#1}\mbox{}}
  \newcommand{\xxxSubParagraphNoStar}[1]{\oldsubparagraph{#1}\mbox{}}
\fi
\makeatother


\providecommand{\tightlist}{%
  \setlength{\itemsep}{0pt}\setlength{\parskip}{0pt}}\usepackage{longtable,booktabs,array}
\usepackage{calc} % for calculating minipage widths
% Correct order of tables after \paragraph or \subparagraph
\usepackage{etoolbox}
\makeatletter
\patchcmd\longtable{\par}{\if@noskipsec\mbox{}\fi\par}{}{}
\makeatother
% Allow footnotes in longtable head/foot
\IfFileExists{footnotehyper.sty}{\usepackage{footnotehyper}}{\usepackage{footnote}}
\makesavenoteenv{longtable}
\usepackage{graphicx}
\makeatletter
\def\maxwidth{\ifdim\Gin@nat@width>\linewidth\linewidth\else\Gin@nat@width\fi}
\def\maxheight{\ifdim\Gin@nat@height>\textheight\textheight\else\Gin@nat@height\fi}
\makeatother
% Scale images if necessary, so that they will not overflow the page
% margins by default, and it is still possible to overwrite the defaults
% using explicit options in \includegraphics[width, height, ...]{}
\setkeys{Gin}{width=\maxwidth,height=\maxheight,keepaspectratio}
% Set default figure placement to htbp
\makeatletter
\def\fps@figure{htbp}
\makeatother

\addtolength{\oddsidemargin}{-.5in}%
\addtolength{\evensidemargin}{-1in}%
\addtolength{\textwidth}{1in}%
\addtolength{\textheight}{1.7in}%
\addtolength{\topmargin}{-1in}%
\makeatletter
\@ifpackageloaded{caption}{}{\usepackage{caption}}
\AtBeginDocument{%
\ifdefined\contentsname
  \renewcommand*\contentsname{Table of contents}
\else
  \newcommand\contentsname{Table of contents}
\fi
\ifdefined\listfigurename
  \renewcommand*\listfigurename{List of Figures}
\else
  \newcommand\listfigurename{List of Figures}
\fi
\ifdefined\listtablename
  \renewcommand*\listtablename{List of Tables}
\else
  \newcommand\listtablename{List of Tables}
\fi
\ifdefined\figurename
  \renewcommand*\figurename{Figure}
\else
  \newcommand\figurename{Figure}
\fi
\ifdefined\tablename
  \renewcommand*\tablename{Table}
\else
  \newcommand\tablename{Table}
\fi
}
\@ifpackageloaded{float}{}{\usepackage{float}}
\floatstyle{ruled}
\@ifundefined{c@chapter}{\newfloat{codelisting}{h}{lop}}{\newfloat{codelisting}{h}{lop}[chapter]}
\floatname{codelisting}{Listing}
\newcommand*\listoflistings{\listof{codelisting}{List of Listings}}
\makeatother
\makeatletter
\makeatother
\makeatletter
\@ifpackageloaded{caption}{}{\usepackage{caption}}
\@ifpackageloaded{subcaption}{}{\usepackage{subcaption}}
\makeatother

\ifLuaTeX
  \usepackage{selnolig}  % disable illegal ligatures
\fi
\usepackage[]{natbib}
\bibliographystyle{agsm}
\usepackage{bookmark}

\IfFileExists{xurl.sty}{\usepackage{xurl}}{} % add URL line breaks if available
\urlstyle{same} % disable monospaced font for URLs
\hypersetup{
  pdftitle={Multilevel Determinants of Player Market Value in Elite European Football: The Role of Performance, Club Prestige, and Player Characteristics},
  pdfauthor={Sagnik Chakrava; Jianing Zou; Kevin Linares; Namit Shrivastava},
  pdfkeywords={football economics, multilevel modeling, player
valuation, club prestige, UEFA coefficient, sports analytics},
  colorlinks=true,
  linkcolor={blue},
  filecolor={Maroon},
  citecolor={Blue},
  urlcolor={Blue},
  pdfcreator={LaTeX via pandoc}}



\begin{document}


\def\spacingset#1{\renewcommand{\baselinestretch}%
{#1}\small\normalsize} \spacingset{1}


%%%%%%%%%%%%%%%%%%%%%%%%%%%%%%%%%%%%%%%%%%%%%%%%%%%%%%%%%%%%%%%%%%%%%%%%%%%%%%

\date{November 24, 2025}
\title{\bf Multilevel Determinants of Player Market Value in Elite
European Football: The Role of Performance, Club Prestige, and Player
Characteristics}
\author{
Sagnik Chakrava\\
and\\Jianing Zou\\
and\\Kevin Linares\\
and\\Namit Shrivastava\\
University of Maryland\\
}
\maketitle

\bigskip
\bigskip
\begin{abstract}
This study investigates the determinants of player market value in elite
European football using a cross-classified multilevel model applied to
7,023 observations between 2013 and 2024. By disentangling individual
reputation from club context, we identify distinct non-linear trends in
market evolution and biological aging. Our results reveal a ``Boom and
Bust'' market cycle, a strict ``Inverted-U'' career arc peaking at age
27, and a significant 42\% valuation premium for key starters (``Squad
Status''). Furthermore, we find evidence of a ``Diversity Premium,''
where clubs with higher international representation garner higher
player valuations, contrasting with a ``Domestic Discount'' for local
players.
\end{abstract}

\noindent%
{\it Keywords:} football economics, multilevel modeling, player
valuation, club prestige, UEFA coefficient, sports analytics
\vfill

\newpage
\spacingset{1.9} % DON'T change the spacing!


\setstretch{2}
\section{Introduction and Research
Motivation}\label{introduction-and-research-motivation}

The valuation of elite athletes in the European soccer transfer market
has evolved from a subjective scouting exercise into a global financial
industry exceeding €7 billion annually. In this high-stakes environment,
accurate valuation is critical for club solvency and competitive
balance. In this study, we operationalize Market Value using
crowd-sourced estimates from Transfermarkt, a metric widely validated in
sports economics literature as a robust proxy for actual transfer fees
and a leading indicator of player market valuation based on demand.

Foundational research in sports economics has long sought to isolate the
determinants of this value. Frick (2007) established that while ``hard''
performance metrics like goals and assists are primary drivers,
valuation follows a strict ``Inverted-U'' biological trajectory,
typically peaking between ages 27 and 28 before depreciating rapidly.
Beyond on-pitch performance, Franck and Nüesch (2012) empirically
demonstrated the ``Superstar Effect,'' finding that non-performance
factors---such as media popularity and brand visibility---can magnify a
player's value significantly, especially at the top end of the wage
distribution. However, these existing studies often model players as
static units nested within a single club. This approach fails to account
for the dynamic reality of modern football, where players transfer
between teams, accumulating value from both their portable human capital
(e.g., experience, age) and the institutional prestige of their current
employer.

Neglecting this cross-classified structure---where player and club
effects intersect rather than nest---likely misestimates the true
drivers of market value. This research addresses this methodological
limitation by applying a cross-classified multilevel model to a decade
of transfer data (2013--2024). By disentangling individual reputation
from club context, this study aims to answer four critical questions
that remain under-explored in nested models:

\begin{enumerate}
\def\labelenumi{\arabic{enumi}.}
\item
  \textbf{Market Evolution (RQ1):} \emph{How has the transfer market
  evolved annually compared to the 2013 baseline?} Understanding the
  non-linear trajectory of market inflation (the ``Boom and Bust''
  cycles) is essential for distinguishing real player value growth from
  sector-wide inflation.
\item
  \textbf{Biological Career Arc (RQ2):} \emph{How does market value
  change as a player moves through their career?} By modeling the
  non-linear effect of age, we aim to validate Frick's (2007) ``peak
  age'' hypothesis, identifying the precise window where physical peak
  and experience intersect.
\item
  \textbf{Squad Status (RQ3):} \emph{Does playing more minutes than the
  average teammate serve as a primary driver of valuation?} While goals
  capture output, they bias towards attackers. We propose that ``Squad
  Status'' (relative minutes played) is a more universal signal of
  manager trust and player utility across all positions.
\item
  \textbf{Contextual Premiums (RQ4):} \emph{Is there a premium for club
  diversity and a discount for domestic players?} As the labor market
  globalizes, we investigate whether clubs with broader international
  recruitment (Diversity Premium) generate higher valuations, and if
  domestic players face a valuation discount compared to imported
  talent.
\end{enumerate}

\section{Method}\label{method}

\emph{Data Source.} We obtained data from
\href{https://www.transfermarkt.com/}{Transfermarkt}and
\href{https://www.uefa.com/}{Union of European Football Associattion}
(UEFA) official rankings, which are both public available websites. At
the time of this study Transfermarket does not have free API access to
ndefinedtheir data repository, but instead the r package
\href{https://jaseziv.github.io/worldfootballR/articles/load-scraped-data.html}{worldfootballR}
allows us to scrap data from their website. However, there are several
limitations to this package such as time out requests and failure to
capture dynamic java tables on webpages. Fortunately, we found the
necessary data through kaggle
\href{https://www.kaggle.com/datasets/davidcariboo/player-scores?resource=download}{\textbf{Football
Data from Transfermarkt}} \textbf{,} which allowed us to focus on the
modeling aspect of this project following data restructuring and
predictor variable engineering.

\emph{Sample Characteristics.} We use the 10 year averaged UEFA
coefficient rankings\footnote{UEFA coefficient rankings are a
  points-based system used by European football's governing body to rank
  national associations and clubs based on their performance in European
  competitions, which determines qualification spots and seeding for
  tournaments like the Champions League and Europa League.} to select
the top 21 European soccer teams as our first stage sampling. In our
second stage non-random sampling, we select soccer players that played
for a full season in one of the 21 clubs from 2013-2024. This resulted
in 2,215 male players representing 7,168 seasons. However, after
examining missing data, which mostly came from missing information on
citizenship status, we were left with 2,159 players representing 7,023
seasons, a 2.5\% loss of player data (e.g., 2\% loss in season data).
The majority of players across clubs and years were in their mid 20s
(\emph{M}=26.5, \emph{SD=}4.6). In this sample, players on average
played 3.2 (\emph{SD}=2.9) seasons, and if we were to droped players
that had less than two seasons to examine player trajectories we would
loose 62\% of our sample. Figure 1 contains the number of players in
each of the 21 clubs across the years.

\includegraphics{MLM_Determinants_of_player_market_value_draft_files/figure-pdf/unnamed-chunk-3-1.pdf}

\emph{Variable selection and Normalization}. Our outcome of interest is
a player's yearly market valuation which we convert from euros to US
dollars. Market valuation in Euro varies by month and we first took the
yearly average for each player prior to transforming to USD. We log
transform market value to 2024 USD because financial data follows a
pareto distribution rather than a normal distribution. Keeping the
outcome on the raw scale would result in extremely large residuals for
high value players, and small for low value players violating the
assumption of homoscedasticity. Log transforming the outcome compresses
the extreme outliers by pulling the long right tail to try and make the
distribution normal. Therefore, this adjustment allows us to compare
real economic power of a player in 2013 versus 2014 without the noise of
currency fluctuation.\footnote{We handle exchange rate fluctuation by
  using historical exchange rates to convert our values to USD. We also
  adjust market value to 2024 USD because 10 million in 2013 is not the
  same as in 2024 due to inflation, and by keeping the raw values we
  would lose real economic value.} Figure 2 shows that market value has
a lot of variability across years, yet most values are at the lower end
which is common for salary data. When we log transform our outcome, we
bring the lower values up and the higher values down.

\includegraphics{MLM_Determinants_of_player_market_value_draft_files/figure-pdf/unnamed-chunk-4-1.pdf}

A list of variables we used for modeling players' market values are
shown in Table 1 alongside how and why they were centered. THIS TABLE
MIGHT NEED TO GO IN The Appendix.

To ensure meaningful model intercepts and convergence, we applied the
following centering strategies:

\begin{itemize}
\item
  \textbf{Outcome:} \texttt{usd\_revenue\_2024\_log} (Natural log of
  inflation-adjusted market value).
\item
  \textbf{Time:} \texttt{season\_year\_c} (Min-centered, 2013=0).
\item
  \textbf{Player Traits:} \texttt{age\_c} (Grand-mean centered) and
  \texttt{height\_c} (Grand-mean centered).
\item
  \textbf{Squad Status:} \texttt{minutes\_played} was
  \textbf{standardized (Z-score)} relative to the specific club-year.
  This isolates a player's status as a ``Starter'' or ``Bench Player''
  relative to their current teammates, removing the noise of team-level
  rotation policies.
\item
  \textbf{Club Context:} \texttt{foreign\_players} (Proportion of
  foreign players, grand-mean centered).
\item
  \textbf{Categorical Controls:} Region, Position, and Citizenship
  (Binary: 1 = Citizen of Club Country).
\end{itemize}

\begin{longtable}[]{@{}
  >{\raggedright\arraybackslash}p{(\columnwidth - 6\tabcolsep) * \real{0.2466}}
  >{\raggedright\arraybackslash}p{(\columnwidth - 6\tabcolsep) * \real{0.2466}}
  >{\raggedright\arraybackslash}p{(\columnwidth - 6\tabcolsep) * \real{0.2466}}
  >{\raggedright\arraybackslash}p{(\columnwidth - 6\tabcolsep) * \real{0.2603}}@{}}
\caption{Table 1. Data Structure for fixed effects, except year as it is
a random slope}\tabularnewline
\toprule\noalign{}
\begin{minipage}[b]{\linewidth}\raggedright
\textbf{Variable Type}
\end{minipage} & \begin{minipage}[b]{\linewidth}\raggedright
\textbf{Variable}
\end{minipage} & \begin{minipage}[b]{\linewidth}\raggedright
Centering Strategy/ Form Factor
\end{minipage} & \begin{minipage}[b]{\linewidth}\raggedright
Justification
\end{minipage} \\
\midrule\noalign{}
\endfirsthead
\toprule\noalign{}
\begin{minipage}[b]{\linewidth}\raggedright
\textbf{Variable Type}
\end{minipage} & \begin{minipage}[b]{\linewidth}\raggedright
\textbf{Variable}
\end{minipage} & \begin{minipage}[b]{\linewidth}\raggedright
Centering Strategy/ Form Factor
\end{minipage} & \begin{minipage}[b]{\linewidth}\raggedright
Justification
\end{minipage} \\
\midrule\noalign{}
\endhead
\bottomrule\noalign{}
\endlastfoot
\textbf{Time} & Season Year & Min-Centering (+ Quadratic) & Establishes
a meaningful starting intercept (2013 =0) \\
\textbf{Contextual, Club} & UEFA Coefficient & Year-Mean Centering &
Isolates relative prestige; removes historical inflation \\
& Foreign Players Proportion & Grand-Mean Centering & Moves baseline to
``Average Club'' (since 0\% is unlikely) \\
\textbf{Player Characteristics} & Age & Grand-Mean (+ Quadratic) &
Captures absolute career stage (Prospect \(\to\) Prime \(\to\)
Veteran) \\
& Height (in cm) & Grand-Mean & Captures premium for absolute physical
size \\
& Position & Reference group, attacker & Attackers tend to be valued
more on average \\
\textbf{Player Demographic} & Citizen of Club Country & None (Raw binary
0/1) & Keeps ``Non-Citizen'' as clear reference group. Separates
individual status from club composition (\texttt{foreign\_players}) \\
& Region of origin & Reference group, Europe & Control variable for
where player is from \\
\textbf{Player Performance} & Goals & Position-Centering & Preserves
natural 0 and marginal value of 1 goal interpretation \\
& Assists & Position-Centering & Preserves natural 0 and marginal value
of 1 assist interpretation \\
& Red Card & raw season count & Meaningful 0 interpretation \\
& Yellow Card & raw season count & Meaningful 0 interpretation \\
& Average minutes played in season & Year-Mean Centering, z-score &
Isolates whether player is a starter or bench player \\
\end{longtable}

\emph{Data Assumptions:} We assume that variables such as height and
positions are time invariant for the duration of this study. At the club
level, we assume mutually exclusiveness in that players can only belong
to one club in a given year. Although a player could have switched clubs
mid-season, we keep this constant for model simplicity. We also assume
that the relationship between age and log-market value is non-linear and
parabolic so we introduce a quadratic term on age. Because we excluded
observations with missing data (listwise deletion), we assume the data
are Missing Completely at Random (MCAR) or that the proportion of
missing cases (\textasciitilde2\%) is negligible and does not introduce
systematic bias.

\hfill\break

\section{Modeling Approach}\label{modeling-approach}

We employ a cross-classified multilevel linear model to partition
variance across observations (Level 1), players (Level 2a), and clubs
(Level 2b). This structure explicitily accounts for players transferring
between clubs and thus treating player and club as cross random factors
rather than a strict three level hierarchy. We found that in this
sample, almost 20\% of players had been in two to seven (17\% in two
clubs) clubs which was the basis for fitting a cross-classified model to
these data. Figure 3 is a visual of a cross-classified model example of
three players that either move through clubs or remain in the same club
throughout their career. Players can change teams over their span of
their careers meaning that the player context can change. They may be a
star player in one team and a bench player on a different team that they
switch to. In a cross-classified model, the player is the constant unit
(i), but the club (j) changes over time (t). If player (I) plays for
club (j1) in year (t0) but switches to club (j2) in year (t1) the fixed
effect in the model, such as minutes played, updates the players status,
as opposed to having static or average predictors. Club attributes, such
as foreign players and UEFA scores, ensures that when player move clubs,
the model swaps the values of \(X_j1\), for \(X_j2\). Player's
characteristics, for instance age and height, move with the player
across time.

\begin{figure}

\centering{

\includegraphics{MLM_Determinants_of_player_market_value_draft_files/figure-pdf/fig-cross-class-1.pdf}

}

\caption{\label{fig-cross-class}Figure 3. Cross-Classified Data
Structure: Players move between clubs over time.}

\end{figure}%

\textbf{Jianing}, can you write this result? This section should be for
the slope, but remember that the slope is on the club not player.

\includegraphics{MLM_Determinants_of_player_market_value_draft_files/figure-pdf/unnamed-chunk-5-1.pdf}

We plot age by yearly market value in 2024 USD (on the log scale) in
Figure 5. for each team and observe almost an upside down u shape. As
players get older you are valued less in the market, same for younger
players.

\includegraphics{MLM_Determinants_of_player_market_value_draft_files/figure-pdf/unnamed-chunk-6-1.pdf}

\hfill\break

\subsection{Model Specification}\label{model-specification}

To formally describe our analysis, we utilize a hierarchical notation
system that decomposes the model into observation, player, and club
levels. This analysis clarifies how the cross-classified structure
(\(i\) players crossed with \(j\) clubs) contributes to the variance in
market value.

Level 1: Observation Model (Time-Varying Effects)

At the observation level (time \(t\), player \(i\), club \(j\)), the
log-transformed market value is a function of a player-club specific
intercept (\(\pi_{0ij}\)), temporal trends, biological aging, and
dynamic performance metrics:

\[
\begin{aligned}
\log(\text{Revenue})_{tij} &= \pi_{0ij} \\
&+ \pi_{1j}(\text{Year}_{t}) + \pi_{2}(\text{Year}^2_{t}) + \pi_{3}(\text{Year}^3_{t}) \\
&+ \pi_{4}(\text{Age}_{it}) + \pi_{5}(\text{Age}_{it})^2 \\
&+ \pi_{6}(\text{Mins}_{Z,it}) + \pi_{7}(\text{Goals}^*_{it}) + \pi_{8}(\text{Assists}^*_{it}) \\
&+ \pi_{9}(\text{Red}_{it}) + \pi_{10}(\text{Yellow}_{it}) \\
&+ \epsilon_{tij}
\end{aligned}
\]

where \(\epsilon_{tij} \sim N(0, \sigma^2_{\epsilon})\).

Level 2a: Player Model (Time-Invariant Traits)

The player-club intercept (\(\pi_{0ij}\)) is determined by static
physical and demographic characteristics, citizenship status relative to
the club, and a player-specific random effect:

\[
\pi_{0ij} = \beta_{0j} + \gamma_{1}(\text{Height}_{i}) + \gamma_{2}(\text{Citizen}_{ij}) + \sum_{p} \gamma_{p} \mathbb{1}_{\text{Position}_i=p} + \sum_{r} \gamma_{r} \mathbb{1}_{\text{Region}_i=r} + u_{0i}
\]

where \(u_{0i} \sim N(0, \sigma^2_{u})\).

Level 2b: Club Model (Contextual Effects)

The club-level baseline (\(\beta_{0j}\)) and the linear time slope
(\(\pi_{1j}\)) are modeled as functions of club internationalization and
club-specific random effects:

\[
\begin{aligned}
\beta_{0j} &= \delta_{00} + \delta_{01}(\text{Foreign}_{tj}) + v_{0j} \\
\pi_{1j} &= \delta_{10} + v_{1j} \\
\pi_{2} &= \delta_{20} \quad \text{(fixed across clubs)} \\
\pi_{3} &= \delta_{30} \quad \text{(fixed across clubs)}
\end{aligned}
\]

where
\(\begin{pmatrix} v_{0j} \\ v_{1j} \end{pmatrix} \sim N \left( \mathbf{0}, \boldsymbol{\Sigma}_{club} \right)\),
with
\(\boldsymbol{\Sigma}_{club} = \begin{pmatrix} \sigma^2_{v_0} & \sigma_{v_0v_1} \\ \sigma_{v_0v_1} & \sigma^2_{v_1} \end{pmatrix}\).

Interpretation of Cross-Classification:

Player effects (\(u_{0i}\)): Capture persistent individual heterogeneity
(e.g., reputation, marketability) that follows a player across clubs.

Club effects (\(v_{0j}, v_{1j}\)): Capture the baseline prestige and
temporal trajectory of the club itself, which affects all players on the
roster.

Cross-classified structure: A player's market value at time \(t\)
depends on both their intrinsic ``brand'' (\(u_{0i}\)) and the club they
are currently with (\(v_{0j} + v_{1j} \cdot \text{Year}_t\)).

\hfill\break

\section{Results}\label{results}

\textbf{Discuss things like REML, LRT test, and mention the models we
fit and how we compared them to arrive at our final model.}

\textbf{Expand on the results a bit more, and double check the
estimates.}

Market Evolution (Q1)

The model reveals a significant cubic trend (p\textless.001) in market
evolution. As shown in Figure 6A, the market followed an ``S-shaped''
trajectory: stagnation (2013--2014), a rapid inflationary ``boom''
(2015--2019), and a post-2020 plateau. This confirms that football
market inflation is cyclical rather than constant.

Biological Career Arc (Q2)

We observe a strong quadratic effect for Age (p\textless.001). Figure 6B
illustrates the ``Inverted-U'' shape, where value peaks during a
player's prime (approx. age 27) and declines significantly in the
veteran stage.

Squad Status Premium (Q3)

Squad status is a massive driver of valuation. Because minutes\_played
was standardized, the coefficient (\(\pi_6 \approx = .35\)) indicates
that a one standard deviation increase in playing time is associated
with a 42\% increase in market value (\(e^{.35} - 1 \approx .42\)). This
confirms that being a ``Starter'' is financially distinct from being a
Rotation Player.

Contextual Effects (Q4-Q5)

We find a ``Diversity Premium'' (\(\delta_{01} = .20\). A 10 percentage
point increase in a club's foreign player proportion increases average
player value by 2\%. Conversely, we find a ``Domestic Discount'' where
homegrown citizens are valued 17\% lower than foreign imports
(\(e^{-.19} - 1 \approx -.17\)).

\begin{figure}

\centering{

\includegraphics{MLM_Determinants_of_player_market_value_draft_files/figure-pdf/fig-results-1.pdf}

}

\caption{\label{fig-results}Key Drivers of Market Value. (A) Cubic
Market Trend. (B) Quadratic Career Arc. (C) Squad Status Premium. (D)
Diversity Premium.}

\end{figure}%

\includegraphics{MLM_Determinants_of_player_market_value_draft_files/figure-pdf/unnamed-chunk-7-1.pdf}

\section{Club-Level Heterogeneity}\label{club-level-heterogeneity}

Figure X displays the estimated club random effects from our
cross-classified model. Substantial heterogeneity exists in both
baseline valuations and temporal trajectories.

\subsection{Baseline Premiums and
Discounts}\label{baseline-premiums-and-discounts}

Manchester United exhibits the largest effect: controlling for age,
position, and performance, United players are valued 68\% higher than
league average (\(\beta_{0j}\) = 0.52, 95\% CI {[}0.31, 0.76{]}).
Barcelona (+56\%, 95\% CI {[}29\%, 93\%{]}), Paris Saint-Germain (+26\%,
95\% CI {[}6\%, 53\%{]}), Chelsea (+25\%, 95\% CI {[}7\%, 50\%{]}),
Manchester City (+24\%, 95\% CI {[}5\%, 49\%{]}), and Bayern Munich
(+22\%, 95\% CI {[}4\%, 48\%{]}) also show significant premiums.

Conversely, seven clubs show significant discounts. Sevilla faces the
steepest penalty: players are valued 28\% below league average
(\(\beta_{0j}\) = -0.33, 95\% CI {[}-0.42, -0.23{]}). Porto (-27\%),
Benfica (-22\%), Roma (-21\%), Napoli (-17\%), Villarreal (-16\%), and
Tottenham (-14\%) also show significant negative effects.

\subsection{Diverging Temporal
Trajectories}\label{diverging-temporal-trajectories}

The random slopes reveal that club fortunes diverged over the study
period. Three clubs showed significantly faster appreciation than league
average: Arsenal (+2.9\%/year, 95\% CI {[}0.8\%, 5.0\%{]}), Liverpool
(+2.7\%/year, 95\% CI {[}0.7\%, 4.8\%{]}), and Tottenham (+2.8\%/year,
95\% CI {[}0.8\%, 4.8\%{]}). These trajectories coincide with improved
sporting performance and, for Liverpool, a Champions League victory in
2019.

Four clubs experienced significantly slower growth (or faster decline):
Manchester United (-2.8\%/year, 95\% CI {[}-4.5\%, -1.0\%{]}), Borussia
Dortmund (-2.3\%/year, 95\% CI {[}-4.1\%, -0.4\%{]}), Juventus
(-2.3\%/year, 95\% CI {[}-4.1\%, -0.5\%{]}), and Barcelona (-1.8\%/year,
95\% CI {[}-3.5\%, -0.1\%{]}). Despite high baseline premiums, these
clubs saw their market advantages erode over time.

\includegraphics{MLM_Determinants_of_player_market_value_draft_files/figure-pdf/unnamed-chunk-8-1.pdf}

\hfill\break

\section{Findings}\label{findings}

This study demonstrates that player market value is driven primarily by
individual reputation and squad status, but is significantly moderated
by club context and market cycles. Our cross-classified approach reveals
that while ``Incubator Clubs'' can accelerate value growth, the market
overall has entered a deceleration phase post-2020. Future research
should address the potential endogeneity of playing time using
instrumental variable approaches.

While our cross-classified model demonstrates strong predictive power
(\(R^2 \approx .89\)) , several limitations warrant consideration.
First, although the Likelihood Ratio Test supported the inclusion of
random slopes for club trajectories, the variance estimates should be
interpreted with caution given the relatively small number of unique
clubs (N=21). Second, diagnostic plots indicated deviations from
normality in the extreme tails of the residual distribution; while our
large sample size (N=7,023) ensures robust estimation of fixed effects,
future iterations could employ robust standard errors or a Gamma GLMM to
better address these extremes. Finally, predictors such as
\texttt{minutes\_played} are likely endogenous---since better players
naturally accrue more playing time---and thus should be interpreted as
associations reflecting ``Squad Status'' rather than purely causal
effects.

Ultimately, this research highlights the necessity of moving beyond
static, single-level analyses to capture the true complexity of the
football transfer market. By rigorously modeling the cross-classified
intersection of individual player trajectories and dynamic club
contexts, we provide a more accurate lens for valuing talent in a
hyper-inflated economy. As the industry continues to professionalize,
such multidimensional frameworks---which account for the interplay of
human capital, institutional prestige, and market cycles---will remain
indispensable for clubs seeking to distinguish genuine value from market
noise.

\newpage

\section{References}\label{references}

Frick, B. (2007). The Football Players' Labor Market: Empirical Evidence
from the Major European Leagues. Scottish Journal of Political Economy,
54(3), 422--446.

Franck, E., \& Nüesch, S. (2012). Talent and/or Popularity: What Does It
Take to Be a Superstar? Economic Inquiry, 50(1), 202--216.

Rosen, S. (1981). The Economics of Superstars. The American Economic
Review, 71(5), 845--858.




\end{document}
